\begin{otherlanguage}{english}
\chapter{Abstract}

A transfusion system cannot manufacture what it distributes: blood, plasma
and platelets depend entirely on donors coming forward voluntarily.
Knowing who donates, how often and for how many years is therefore an
operational prerequisite for planning collection.
This thesis addresses the question with real data:
the donation records of the \textbf{Interhospital Department of Transfusion
Medicine (DIMT)} of Trieste and Gorizia, two provinces of Friuli Venezia
Giulia whose population is older on average than the rest of Italy and whose
donor associations are particularly well established.

The work starts from 27 Excel files with inconsistent schemas and naming
conventions and, through a Python pipeline documented at every stage, merges
them into a single analytical dataset of \textbf{353\,671 donation events}
from \textbf{43\,299 distinct donors} over \textbf{2009--2025}. Four analyses
are built on it: a descriptive exploration of the donor population and its
trends over time, related for the first time to the resident population; a
retention study by entry cohort; a projection of donation volume to 2031,
selected through a rolling-origin backtest among seventeen candidate models; and
a count regression model of individual donation frequency. The results are
also made explorable in an interactive R Shiny dashboard.

The picture is one of slow, steady contraction. Annual volume falls by
\textbf{28.1\%} between 2009 and 2025, with most of the decline in the first
half of the period before it levels off; new entries fall even more; and
individual intensity --- donations per active donor --- stays between $1.84$
and $1.97$ across all seventeen years. The system, in other words, did not
intensify to make up for lower turnover: it shrank and then settled on a
smaller base, helped in part by falling demand.

Relating donations to the resident population of Trieste and Gorizia separates
demography from behaviour, and the result is not what one would expect. The
decline in the population of eligible age explains only \textbf{12.2\%} of the
contraction; nearly all the rest is a genuine loss of participation, from
$106.9$ to $80.1$ donations per thousand residents. Nor is the loss where one
would look for it: relative to population, \textbf{people aged 20--24 have the
highest donation rate in every year observed}, while the 35--44 age groups
lose between 36 and 41 per cent. The problem with young donors is not that
they fail to start --- they account for half of all entries --- but that they
do not come back: \textbf{40.6\%} never return after their first donation,
against $33.3\%$ of those who start later, and the decreasing hazard function
points to a threshold around the third to fifth year at which the habit
consolidates.

The frequency model shows that the dominant effects are regulatory rather than
motivational (sex, IRR $0.585$; plasmapheresis, IRR $1.77$: different minimum
intervals and annual caps) and that, once years of donor career are controlled
for, \textbf{age has no effect at all}: what sets a frequent donor apart is how
long they have been donating, not how old they are. Consistently, the
projections point to about \textbf{97\,200 donations} in total over 2026--2031: a first
year clearly below 2025, then a contraction of just under one percentage point
a year, with no expected shift in the mix of donation types. A slow thinning,
not a break.

\end{otherlanguage}
